% kernel/engineering_context-DE.tex

\chapter{Anwendbarkeit aus verteilten Qualifikationen}
\label{chap:engineering-context}

Der Kandidat besitzt Objektbezug, konstruktive Beschreibung und einen
nachvollziehbaren Vorgänger. Nun fragt das Projekt, ob er für eine beabsichtigte
Betriebsgeschwindigkeit genügt. Diese Frage lässt sich nicht durch ein
allgemeines Kontextetikett beantworten.

Der aktive Knowledge Kernel etabliert kein einheitliches Objekt ``Engineering
Context''. Die Thesis hält die erforderlichen Qualifikationen deshalb bei den
Konzepten, denen sie gehören.

\section{Der Zweck bestimmt die erforderliche Unterstützung}

Der angegebene Zweck ist die Beurteilung des Kandidaten für den Betrieb mit
einer bestimmten Geschwindigkeit, nicht nur seine Darstellung. Dafür werden
anwendbare Krümmungs- und Überhöhungsgrenzen, Fahrzeug- und Betriebsannahmen,
Toleranzen und eine geeignete Evidenzqualität benötigt. Eine CAD-Vorschau kann
für die visuelle Prüfung genügen und für die betriebliche Beurteilung
unzureichend sein. Der Zweck verändert die erforderliche Unterstützung, nicht
die Identität des Alignments.

\section{Realisierung verleiht der Konstruktion metrische Bedeutung}

Die intrinsische Folge und $\kappa(s)$ hängen nicht von einer projizierten
Zeichnung ab. Eine Metric Realization interpretiert sie als messbare Geometrie
unter einem Realization Context: metrischer Raum, Referenzsystem oder
Oberfläche, Frames, Konventionen, Genauigkeit und
Anwendbarkeitsabhängigkeiten. Das Koordinatenreferenzsystem ist eine
Abhängigkeit dieser Beschreibung, nicht ihr Owner und nicht die Identität des
Alignments.

Auch die Stationierung zur Adressierung des betroffenen Intervalls muss
ausdrücklich auf die intrinsische Längsparametrisierung bezogen werden. Eine
Kilometrierungskonvention oder Stationsgleichung kann die Adressierung einer
Position ändern, ohne die konstruktive Folge zu verändern. Verwenden eine
erdbezogene und eine lokale planare Beurteilung verschiedene gültige
Realisierungen, müssen ihre Ergebnisse diese Kontexte behalten, statt
stillschweigend verbunden zu werden.

\section{Weitere Qualifikationen bleiben verteilt}

Berechnung und Quelle tragen zeitliche Gültigkeit und Provenienz. Beabsichtigte
Geschwindigkeit und Bewertungszweck tragen ihren eigenen Geltungsbereich.
Anwendbare Normen, Projektregeln und Überhöhungsbedingungen behalten ihre Owner
und Gültigkeitsbereiche. Projektphase und Workflow-Status können die betrachtete
Aussage qualifizieren, bilden aber keine abgeschlossene Kernel-Abstraktion und
werden hier nicht zu einer solchen erhoben.

Der Kandidat ist nur anwendbar, wenn diese Abhängigkeiten gemeinsam die
angegebene Konsequenz tragen. Anwendbarkeit ist ein Ergebnis dieser
qualifizierten Beziehung, keine durch erfolgreiche Berechnung verliehene
Eigenschaft.

\section{Die erste operationale Prüfung}

Das Projekt kann nun präzise fragen: Erfüllt der vorgeschlagene Krümmungs- und
Überhöhungszustand unter der gewählten Metric Realization und der angegebenen
Geschwindigkeit die anwendbaren Bedingungen mit Evidenz der erforderlichen
Qualität und Gültigkeit? Ein negatives oder unbestimmtes Ergebnis lässt das
mathematische Resultat berechnet, aber nicht anwendbar. Ein positives Ergebnis
macht es für diesen Zweck nutzbar, aber noch nicht autoritativ.

Damit diese Prüfung Quelle und erzeugte Kurve nicht verwechselt, trennen die
nächsten Kapitel Repräsentationen, Beobachtungen und angenommenes Wissen.
